\DocumentMetadata{
  pdfversion=2.0,
  pdfstandard=ua-2,
  lang=en-US,
  tagging=on,
  tagging-setup={math/setup=mathml-SE,tabsorder=structure}
}
\documentclass[11pt,letterpaper]{article}
\usepackage[margin=0.68in,includefoot]{geometry}
\usepackage{newcomputermodern}
\usepackage{amsmath,amscd,mathtools}
\usepackage{tikz,adjustbox}
\usetikzlibrary{arrows.meta,decorations.markings,positioning}
\providecommand{\square}{\mdlgwhtsquare}
\usepackage[hidelinks]{hyperref}
\hypersetup{
  pdftitle={Topology PhD Exam, May 2025},
  pdfauthor={University of Florida Department of Mathematics},
  pdfsubject={Graduate mathematics examination},
  pdfdisplaydoctitle=true,
  bookmarksopen=true
}
\setcounter{secnumdepth}{0}
\setlength{\parindent}{0pt}
\setlength{\parskip}{0.28em}
\setlength{\emergencystretch}{3em}
\setlength{\leftmargini}{2.15em}
\setlength{\leftmarginii}{2.65em}
\setlength{\leftmarginiii}{3em}
\renewcommand{\labelenumi}{\textbf{\theenumi.}}
\pagestyle{empty}
\raggedbottom
\allowdisplaybreaks
\newenvironment{examproblems}[1]{%
  \begin{enumerate}%
  \setcounter{enumi}{#1}%
  \setlength{\topsep}{0.35em}%
  \setlength{\partopsep}{0pt}%
  \setlength{\itemsep}{0.48em}%
  \setlength{\parsep}{0.18em}%
}{\end{enumerate}}
\newenvironment{examsubparts}{%
  \begin{enumerate}%
  \setlength{\topsep}{0.18em}%
  \setlength{\partopsep}{0pt}%
  \setlength{\itemsep}{0.16em}%
  \setlength{\parsep}{0.1em}%
}{\end{enumerate}}
\newenvironment{examinstructions}{%
  \par\begingroup\itshape
}{%
  \par\endgroup\vspace{0.18em}
}
\makeatletter
\renewcommand\section{\@startsection{section}{1}{0pt}%
  {-1.25ex plus -.25ex minus -.1ex}%
  {0.55ex plus .1ex}%
  {\normalfont\large\bfseries}}
\renewcommand{\@maketitle}{%
  \newpage\null\vskip -1.2em%
  {\footnotesize\bfseries
    \href{https://www.ufl.edu/}{University of Florida}, \href{https://math.ufl.edu/}{Department of Mathematics}\par}%
  \vskip 0.18em%
  \tagstructbegin{tag=Title}%
    {\LARGE\bfseries\href{https://gma.math.ufl.edu/exams/topology-phd/}{Topology PhD Exam}, May 2025\par}%
  \tagstructend%
  \vskip 0.35em%
  \tagmcbegin{artifact}%
    \hrule height 1pt%
  \tagmcend%
  \vskip 0.55em%
}
\makeatother
\title{Topology PhD Exam, May 2025}
\author{}
\date{}
\begin{document}
\maketitle
\thispagestyle{empty}
\begin{examinstructions}
Be neat and do not write too small.
\end{examinstructions}
\begin{examinstructions}
For the first five problems, show all of your work and support all statements.
\end{examinstructions}
\begin{examproblems}{0}
\item
Show that the product of two connected spaces is connected.
\item
Prove that every metrizable space \((X,d)\) is normal.
\item
Let~\(X\) be the~\(\Delta\)-complex drawn below, which has two 0-simplices~\(v\) and~\(w\), four 1-simplices \(a,b,c,d\), and three 2-simplices \(S,T,U\). Compute the simplicial homology \(H_i(X)\) for \(i=0,1,2\).
\par\smallskip
\begin{center}
\begin{adjustbox}{max width=.8\linewidth,max totalheight=10\baselineskip}
\begin{tikzpicture}[
  alt={Triangular diagram illustrating edge identifications and orientations. All three corner vertices are labeled v, with an interior vertex w. Edges from w to the corners divide the triangle into regions S, T, and U. The three boundary edges are all labeled a and oriented cyclically, each with one arrowhead. The edge from the top vertex to w is labeled b and has two arrowheads. The lower-left edge to w is labeled c and has three arrowheads. The lower-right edge to w is labeled d and has four arrowheads.},
  scale=.5
]
\tikzset{
  edge/.style={
    draw=black,
    line width=1.2pt,
    line cap=round,
    line join=round
  },
  vertex/.style={
    circle,
    fill=black,
    inner sep=0pt,
    minimum size=2mm
  },
  % #1 = direction: 1 follows the path, -1 reverses it
  % #2 = number of packed arrowheads
  feather mark/.style 2 args={
    postaction={decorate},
    decoration={
      markings,
      mark=at position .5 with {
        \begin{scope}[xscale=#1]
          \foreach \i in {1,...,#2}{
            % Spacing between arrowheads.
            \pgfmathsetmacro{\dx}{(\i-(#2+1)/2)*3.5}
            \draw[
              black,
              line width=1.1pt,
              line cap=round,
              line join=round,
              xshift=\dx pt
            ]
              (-4pt,-3pt) -- (0pt,0pt) -- (-4pt,3pt);
          }
        \end{scope}
      }
    }
  }
}

\coordinate (Top)    at ( 0.00,5.45);
\coordinate (Left)   at (-3.20,0.00);
\coordinate (Right)  at ( 3.20,0.00);
\coordinate (Center) at ( 0.00,2.10);

% Shaded outer triangle.
\fill[black!20] (Top) -- (Left) -- (Right) -- cycle;

% Outer edges.
\draw[edge,feather mark={1}{1}]
  (Top) -- node[pos=.5,above left=0pt] {$a$} (Left);

\draw[edge,feather mark={1}{1}]
  (Left) -- node[pos=.5,below right=0pt] {$a$} (Right);

\draw[edge,feather mark={1}{1}]
  (Right) -- node[pos=.5,above right=0pt] {$a$} (Top);

% Inner edges.
\draw[edge,feather mark={1}{2}]
  (Top) -- node[pos=.5,below right=0pt] {$b$} (Center);

\draw[edge,feather mark={1}{3}]
  (Left) -- node[pos=.5,below right=0pt] {$c$} (Center);

\draw[edge,feather mark={1}{4}]
  (Right) -- node[pos=.5,below left=0pt] {$d$} (Center);

% Region labels.
\node at (-1,2.4) {$S$};
\node at ( 1,2.4) {$T$};
\node at ( 0,0.7) {$U$};

% Vertices.
\node[vertex,label=above:$v$] at (Top) {};
\node[vertex,label=below left:$v$] at (Left) {};
\node[vertex,label=below right:$v$] at (Right) {};
\node[vertex,label=below:$w$] at (Center) {};
\end{tikzpicture}
\end{adjustbox}
\end{center}
\smallskip
\item
This problem has three parts.
\begin{examsubparts}
\item[\textnormal{(a)}]
For \(0\to A\to B\to C\to 0\) a short exact sequence of finitely generated abelian groups, prove \(\operatorname{rank}(B)=\operatorname{rank}(A)+\operatorname{rank}(C)\).
\item[\textnormal{(b)}]
For \(\cdots\to C_n\xrightarrow{d_n}C_{n-1}\to\cdots\) a chain complex of finitely generated abelian groups, prove 
\[
\operatorname{rank}(C_n)=\operatorname{rank}(\ker d_n)+\operatorname{rank}(\operatorname{im}d_n)
\]
 and 
\[
\operatorname{rank}(\ker d_n)=\operatorname{rank}(\operatorname{im}d_{n+1})+\operatorname{rank}\left(\frac{\ker d_n}{\operatorname{im}d_{n+1}}\right).
\]
\item[\textnormal{(c)}]
Prove a finite CW complex~\(X\) with \(c_n\) cells in dimension~\(n\) satisfies 
\[
\sum_n(-1)^n c_n=\sum_n(-1)^n\operatorname{rank}(H_n(X)).
\]
\end{examsubparts}
\item
This problem has two parts.
\begin{examsubparts}
\item[\textnormal{(a)}]
Let \(n>m\). Show there are no maps from \(\mathbb{RP}^n\) to \(\mathbb{RP}^m\) inducing a nontrivial map \(H^1(\mathbb{RP}^m;\mathbb Z_2)\to H^1(\mathbb{RP}^n;\mathbb Z_2)\).
\item[\textnormal{(b)}]
Prove the Borsuk–Ulam theorem as follows. Suppose on the contrary that \(f:S^n\to\mathbb R^n\) satisfies \(f(x)\ne f(-x)\) for all~\(x\). Define \(g:S^n\to S^{n-1}\) by 
\[
g(x)=\frac{f(x)-f(-x)}{\lVert f(x)-f(-x)\rVert},
\]
 so \(g(-x)=-g(x)\) and~\(g\) induces a map \(\mathbb{RP}^n\to\mathbb{RP}^{n-1}\). Show that (a) applies to this map.
\end{examsubparts}
\end{examproblems}
\begin{examinstructions}
Answer the following ten problems with complete definitions, complete statements, an example, or a short proof.
\end{examinstructions}
\begin{examproblems}{5}
\item
State the Urysohn Lemma.
\item
Give an example of a space that is path-connected but not locally path-connected.
\item
Use the Seifert–van Kampen theorem to derive the fundamental group of \(M_g\), the orientable surface of genus~\(g\), for \(g\ge 1\).
\item
Show that for any connected closed orientable~\(n\)-manifold~\(M\) there is a degree~\(1\) map \(M\to S^n\).
\item
Show that every map \(f:\mathbb{CP}^n\to\mathbb{CP}^n\) has a fixed point if~\(n\) is even.
\item
Show that \(S^1\vee S^1\) admits a covering space whose fundamental group is nontrivial and abelian.
\item
Let \(S^1\times D^2\) be the solid torus, and let \(A\subseteq S^1\times D^2\) be the circle shown in the figure. Prove that no retraction \(r:S^1\times D^2\to A\) exists.
\par\smallskip
\begin{center}
\begin{adjustbox}{max width=.8\linewidth,max totalheight=6\baselineskip}
\begin{tikzpicture}[
  alt={Diagram of a thick, coiled ribbon forming two concentric loops inside a thin oval boundary. A vertical loop rises from the top center, and the inner coil curls inward near the bottom. At lower left, a curved arrow labeled A points to the outer boundary.},
  x=0.01cm,
  y=-0.01cm
]
% Preserve the proportions and whitespace of the reference image.
\path[use as bounding box] (0,0) rectangle (903,480);

\tikzset{
  heavy/.style={
    draw=black,
    line width=2.6pt,
    line cap=butt,
    line join=round
  },
  fine/.style={
    draw=black,
    line width=1.15pt,
    line cap=butt,
    line join=round
  }
}

% Outer thick curve: left segment.
\draw[heavy]
  (476,77)
    .. controls (315,76) and (176,132) .. (176,210)
    .. controls (176,281) and (302,336) .. (443,345);

% Outer thick curve: right segment and lower curl.
\draw[heavy]
  (505,77)
    .. controls (665,78) and (775,132) .. (774,210)
    .. controls (773,279) and (668,327) .. (552,341)
    .. controls (500,348) and (451,339) .. (442,321)
    .. controls (438,312) and (443,303) .. (454,299);

% Inner thick curve: left segment and downward hook.
\draw[heavy]
  (477,130)
    .. controls (349,129) and (262,161) .. (262,210)
    .. controls (262,252) and (348,281) .. (437,284)
    .. controls (471,285) and (494,299) .. (498,314)
    .. controls (499,319) and (498,323) .. (496,325);

% Inner thick curve: right segment.
\draw[heavy]
  (506,130)
    .. controls (615,130) and (689,162) .. (689,208)
    .. controls (689,249) and (617,278) .. (507,280);

% Outer boundary, interrupted at the lower-left arrow.
\draw[fine]
  (204,338)
    .. controls (133,305) and (94,261) .. (94,215)
    .. controls (94,117) and (263,38) .. (470,38)
    .. controls (679,38) and (850,116) .. (850,215)
    .. controls (850,313) and (680,391) .. (472,391)
    .. controls (381,391) and (295,375) .. (229,349);

% Inner visible boundary.
\draw[fine]
  (335,190)
    .. controls (335,221) and (395,246) .. (471,246)
    .. controls (547,246) and (607,221) .. (607,190);

% Top arc of the central donut hole.
\draw[fine]
  (350,213)
    .. controls (390,193) and (430,183) .. (470,183)
    .. controls (510,183) and (552,193) .. (592,213);

% Vertical loop; the left side is interrupted at the crossings.
\draw[fine]
  (455,60)
    .. controls (459,47) and (465,39) .. (470,39)
    .. controls (488,39) and (493,70) .. (493,109)
    .. controls (493,150) and (485,182) .. (470,182);

\draw[fine]
  (470,182)
    .. controls (459,182) and (452,168) .. (449,150);

\draw[fine]
  (450,116)
    .. controls (449,108) and (449,101) .. (450,93);

% Label and curved pointer.
\draw[
  fine,
  -{Stealth[length=2.7mm,width=1.7mm]}
]
  (137,390)
    .. controls (178,382) and (218,344) .. (249,309);

\node[scale=2.7] at (94,389) {$A$};
\end{tikzpicture}
\end{adjustbox}
\end{center}
\smallskip
\item
Show that if \(f:X\to Y\) is a continuous bijection with~\(X\) compact and~\(Y\) Hausdorff, then~\(f\) is a homeomorphism.
\item
State the Baire Category Theorem, and use it to prove that \([0,1]\) is uncountable.
\item
Does there exist a closed orientable 14-dimensional manifold~\(M\) with \(H^7(M;\mathbb Z)\cong\mathbb Z\)?
\end{examproblems}
\end{document}
